%=====================================================================
\section{Decision Orchestration}
\label{sec:method}

\paragraph{Setting.}
An annotation scheme defines a structured annotation $y$ for a sentence or document $x$: in PARSEME, a set of categorized, possibly discontinuous token groups per sentence; in \umr, one rooted graph per sentence plus document-level temporal, modal and coreference relations. In the cold start we have a guideline, the resources it refers to (for \umr, PropBank frames; \citealp{palmer2005propbank}), a scorer and a handful of gold annotations. We build, and later learn, a procedure $\Pi$ that maps $x$ to $y$.

\paragraph{Skills, registry and execution.}
$\Pi$ consists of \textbf{skills}, a \textbf{registry} and numeric parameters $\theta$. A skill has an input schema, an output schema (its \emph{emit}) and an implementation: code (a \emph{rule node}) or a model call (a \emph{decision node}). The registry declares \emph{waves} that run in order; within a wave, a dispatcher maps fields of a skill's emit to the next skill, growing a tree of nodes until every branch ends in a leaf. Control flow is thus a fixed function of typed decisions: answers can change tree shape and call counts, but the model cannot choose an undeclared skill.
In \umr (Table~\ref{tab:nodes}, appendix), \skill{sent\_type} routes a span to \skill{discourse}, \skill{predicate}, \skill{np\_phrase} or \skill{special\_entity}, or stops it as a leaf; fragments leave typed placeholders that go back to \skill{sent\_type}; rule nodes re-dispatch uncovered tokens; later waves merge co-referring nodes and add document-level relations. The PARSEME procedure is the smallest instance: a rule node enumerates candidates from the tokens, lemmas, parts of speech and dependency relations supplied with the data, a closed decision node asks one question per candidate, and the output handler resolves overlaps and writes the \texttt{.cupt} file (Appendix~\ref{app:parseme}).

\paragraph{Two kinds of decision nodes.}
\label{sec:decisions}
A decision node has two implementations that produce the \emph{same} emit, so they can be swapped node by node. In a \textbf{fragment decision}, an LLM receives the skill's instructions (the relevant part of the guideline), the span and the sentence, may call tools for examples and frames, and answers through the emit schema, for example with a graph fragment; it never writes a whole graph.
In a \textbf{closed decision}, the skill's code enumerates candidates from the input (the words of a span, the senses and roles of a verb, the allowed categories) and asks \emph{closed questions}: yes/no, a choice among options that each carry a one-line description from the guideline, or a graded score. A decision model returns a distribution over the options and no text. The skill is a coroutine that yields rounds of questions and returns an emit built by rules, which read the distributions against thresholds $\theta$. \skill{predicate}, for instance, asks for the head, then its sense, aspect, modal strength and polarity, then the role of every content word (Figure~\ref{fig:overview}b); a PARSEME candidate's question lists the categories its kind allows plus \emph{none} (``\emph{gave + up}: VPC.full, VPC.semi, none''), and the candidate is kept when its best option is a category with $p\geq 0.5$. Closed mode needs no model that can write structured text, and every decision comes with a probability.

\paragraph{Traces and attribution.}
\label{sec:trace}
Every node records its input, its decision (in closed mode, the questions and full distributions) and its emit. The scorer's alignment of predicted to gold units lets us record the node that produced each aligned predicted unit; gold units that nothing predicted have no producer, such as a PARSEME expression that no rule proposed. This is output provenance, not causal attribution: upstream routing, candidate omissions and alignment errors can change downstream losses. A fine-grained \umr judge scores 25 aspects separately (Appendix~\ref{app:judge}). Reflection reads these records.
